% GENERATED FILE. Edit canonical lessons, then run npm run generate:books.
% canonical-source-hash: fnv1a64:5088c9d7d97636f0
% canonical-lessons: GE-C365-es-sei-denn, GE-C365-ausser-wenn, GE-C365-andernfalls, GE-C365-ansonsten, GE-C365-notfalls

\chapter{Es sei denn …}
\label{ch:ge-es-sei-denn}
\hlchaptermodality{365}

\begin{chapteropening}
\textbf{By the end of this chapter:} \emph{I can name an exception to a plan and say unless, except when, otherwise and if need be.}

Complete the last lesson of chapter 365: I can name an exception to a plan and say unless, except when, otherwise and if need be.
\end{chapteropening}

\section[es sei denn]{es sei denn — unless}

\label{lesson:GE-C365-es-sei-denn}



\begin{quote}
Before the new one: say the German for provided that, then the German for as far as.
\end{quote}

\subsection*{You'll want to know: es sei denn}
\textbf{es sei denn} — \textquotedblleft{}unless\textquotedblright{}.

\begin{quote}
Wir fahren morgen ans Meer, es sei denn, es regnet.
\end{quote}

\begin{quote}
\emph{We are going to the seaside tomorrow, unless it rains.}
\end{quote}

After \emph{es sei denn} the clause keeps normal order, with the verb second: \emph{es sei denn, es regnet}. \emph{Sei} is the subjunctive (Konjunktiv I) of \emph{sein}.

\begin{cousinweb}
\emph{es sei denn} is a set phrase built on \emph{sei}, the subjunctive of \emph{sein} \textquotedblleft{}to be\textquotedblright{}, and \emph{denn} \textquotedblleft{}then\textquotedblright{}, of the family of \emph{dann} (EtymDB): \textquotedblleft{}unless it be that …\textquotedblright{}.
\end{cousinweb}

\begin{grammarlens}[title={What you've built}]
One more word for telling, explaining and describing at B1.
\end{grammarlens}

\subsection*{Your turn}
\begin{itemize}
\raggedright
  \item \emph{Say it:} \emph{es sei denn}
  \item \emph{Say it:} \emph{es sei denn}, once more
  \item \emph{Say it:} a sentence of your own with \emph{es sei denn}
  \item \emph{Recall:} say the German for provided that, then the German for as far as, then say \emph{es sei denn} again
\end{itemize}

\begin{tcolorbox}[breakable,colback=teal!4,colframe=teal!35!black,title={Before you move on}]
What does es sei denn mean? (\textquotedblleft{}Unless\textquotedblright{}.) Say it once more.
\end{tcolorbox}
\section[außer wenn]{außer wenn — except when, unless}

\label{lesson:GE-C365-ausser-wenn}



\begin{quote}
Before the new one: say the German for as far as, then the German for unless.
\end{quote}

\subsection*{You'll want to know: außer wenn}
\textbf{außer wenn} — \textquotedblleft{}except when, unless\textquotedblright{}.

\begin{quote}
Ich gehe jeden Tag zu Fuß, außer wenn es schneit.
\end{quote}

\begin{quote}
\emph{I walk every day, except when it snows.}
\end{quote}

\emph{Wenn} sends the verb to the end: \emph{außer wenn es schneit}.

\begin{cousinweb}
\emph{außer wenn} joins \emph{außer} \textquotedblleft{}except, outside\textquotedblright{}, from Middle High German \emph{ūzer} and Old High German \emph{ūzar} (EtymDB), to \emph{wenn} \textquotedblleft{}if, when\textquotedblright{}.
\end{cousinweb}

\begin{grammarlens}[title={What you've built}]
One more word for telling, explaining and describing at B1.
\end{grammarlens}

\subsection*{Your turn}
\begin{itemize}
\raggedright
  \item \emph{Say it:} \emph{außer wenn}
  \item \emph{Say it:} \emph{außer wenn}, once more
  \item \emph{Say it:} a sentence of your own with \emph{außer wenn}
  \item \emph{Recall:} say the German for as far as, then the German for unless, then say \emph{außer wenn} again
\end{itemize}

\begin{tcolorbox}[breakable,colback=teal!4,colframe=teal!35!black,title={Before you move on}]
What does außer wenn mean? (\textquotedblleft{}Except when, unless\textquotedblright{}.) Say it once more.
\end{tcolorbox}
\section[andernfalls]{andernfalls — otherwise, if not}

\label{lesson:GE-C365-andernfalls}



\begin{quote}
Before the new one: say the German for unless, then the German for except when.
\end{quote}

\subsection*{You'll want to know: andernfalls}
\textbf{andernfalls} — \textquotedblleft{}otherwise, if not\textquotedblright{}.

\begin{quote}
Wir müssen jetzt gehen, andernfalls verpassen wir den Bus.
\end{quote}

\begin{quote}
\emph{We have to go now; otherwise we will miss the bus.}
\end{quote}

The verb comes right after it: \emph{andernfalls verpassen wir …}. A formal \emph{sonst}.

\begin{cousinweb}
\emph{andernfalls} grew out of \emph{andern Falls} \textquotedblleft{}in the other case\textquotedblright{}, from \emph{ander} \textquotedblleft{}other\textquotedblright{} and \emph{Fall} \textquotedblleft{}fall, case\textquotedblright{}, Old High German \emph{val} (EtymDB), with the adverb ending \emph{-s}.
\end{cousinweb}

\begin{grammarlens}[title={What you've built}]
One more word for telling, explaining and describing at B1.
\end{grammarlens}

\subsection*{Your turn}
\begin{itemize}
\raggedright
  \item \emph{Say it:} \emph{andernfalls}
  \item \emph{Say it:} \emph{andernfalls}, once more
  \item \emph{Say it:} a sentence of your own with \emph{andernfalls}
  \item \emph{Recall:} say the German for unless, then the German for except when, then say \emph{andernfalls} again
\end{itemize}

\begin{tcolorbox}[breakable,colback=teal!4,colframe=teal!35!black,title={Before you move on}]
What does andernfalls mean? (\textquotedblleft{}Otherwise, if not\textquotedblright{}.) Say it once more.
\end{tcolorbox}
\section[ansonsten]{ansonsten — apart from that; otherwise}

\label{lesson:GE-C365-ansonsten}



\begin{quote}
Before the new one: say the German for except when, then the German for otherwise.
\end{quote}

\subsection*{You'll want to know: ansonsten}
\textbf{ansonsten} — \textquotedblleft{}apart from that; otherwise\textquotedblright{}.

\begin{quote}
Die Suppe war etwas salzig, ansonsten war das Essen sehr gut.
\end{quote}

\begin{quote}
\emph{The soup was a bit salty; apart from that, the meal was very good.}
\end{quote}

Two senses: apart from that, and otherwise, like \emph{sonst}. The verb comes right after it.

\begin{cousinweb}
\emph{ansonsten} is built from \emph{an} and \emph{sonst} \textquotedblleft{}otherwise\textquotedblright{}, inherited from Middle High German \emph{sunst} (EtymDB).
\end{cousinweb}

\begin{grammarlens}[title={What you've built}]
One more word for telling, explaining and describing at B1.
\end{grammarlens}

\subsection*{Your turn}
\begin{itemize}
\raggedright
  \item \emph{Say it:} \emph{ansonsten}
  \item \emph{Say it:} \emph{ansonsten}, once more
  \item \emph{Say it:} a sentence of your own with \emph{ansonsten}
  \item \emph{Recall:} say the German for except when, then the German for otherwise, then say \emph{ansonsten} again
\end{itemize}

\begin{tcolorbox}[breakable,colback=teal!4,colframe=teal!35!black,title={Before you move on}]
What does ansonsten mean? (\textquotedblleft{}Apart from that; otherwise\textquotedblright{}.) Say it once more.
\end{tcolorbox}
\section[notfalls]{notfalls — if need be, if necessary}

\label{lesson:GE-C365-notfalls}



\begin{quote}
Before the new one: say the German for otherwise, then the German for apart from that.
\end{quote}

\subsection*{You'll want to know: notfalls}
\textbf{notfalls} — \textquotedblleft{}if need be, if necessary\textquotedblright{}.

\begin{quote}
Notfalls schlafen wir im Auto.
\end{quote}

\begin{quote}
\emph{If need be, we will sleep in the car.}
\end{quote}

The verb comes right after it: \emph{notfalls schlafen wir …}. From \emph{der Notfall} — the emergency.

\begin{cousinweb}
\emph{notfalls} is built from \emph{Not} \textquotedblleft{}need\textquotedblright{}, a cousin of English \emph{need} (EtymDB), and \emph{Fall} \textquotedblleft{}case\textquotedblright{}, Old High German \emph{val} (EtymDB), with the adverb ending \emph{-s}: in case of need.
\end{cousinweb}

\begin{grammarlens}[title={What you've built}]
One more word for telling, explaining and describing at B1.
\end{grammarlens}

\subsection*{Your turn}
\begin{itemize}
\raggedright
  \item \emph{Say it:} \emph{notfalls}
  \item \emph{Say it:} \emph{notfalls}, once more
  \item \emph{Say it:} a sentence of your own with \emph{notfalls}
  \item \emph{Recall:} say the German for otherwise, then the German for apart from that, then say \emph{notfalls} again
\end{itemize}

\begin{tcolorbox}[breakable,colback=teal!4,colframe=teal!35!black,title={Before you move on}]
What does notfalls mean? (\textquotedblleft{}If need be, if necessary\textquotedblright{}.) Say it once more.
\end{tcolorbox}
